\Needspace{6\baselineskip}\Article{345}
\textsc{Theorem.} \textit{Let \((f,\allowbreak \lambda),\allowbreak (f,\allowbreak \mu)\) be two similar periods, identical or distinct,} \textit{and let \((f,\allowbreak \lambda)\) consist of the roots \([\lambda],\allowbreak [\lambda'],\allowbreak [\lambda'']\), etc. Then the product of \((f,\allowbreak \lambda)\) and \((f,\allowbreak \mu)\)} \textit{will be a sum of \(f\) similar periods, namely} \[=(f,\lambda+\mu)+(f,\lambda'+\mu)+(f,\lambda''+\mu)+\mathrm{etc.}=W\]

\textit{Proof.} As above, let \(n-1=ef\); let \(g\) be a primitive root modulo \(n\), and \(h=g^e\), whence, by what precedes, \((f,\allowbreak \lambda)=(f,\allowbreak \lambda h)=(f,\allowbreak \lambda hh)\), etc. The required product will therefore be \[=[\mu].(f,\lambda)+[\mu h].(f,\lambda h)+[\mu hh].(f,\lambda hh)+\mathrm{etc.}\] and hence \[\begin{array}{@{}l@{\;}l@{\;}l@{}}
=[\lambda+\mu]&+[\lambda h+\mu]&\ldots+[\lambda h^{f-1}+\mu]\\
+[\lambda h+\mu h]&+[\lambda hh+\mu h]&\ldots+[\lambda h^f+\mu h]\\
+[\lambda hh+\mu hh]&+[\lambda h^3+\mu hh]&\ldots+[\lambda h^{f+1}+\mu hh]\ \mathrm{etc.}
\end{array}\] This expression contains altogether \(ff\) roots. If each vertical column is now summed separately, we clearly obtain \[(f,\lambda+\mu)+(f,\lambda h+\mu)+\ldots+(f,\lambda h^{f-1}+\mu)\] This expression is readily seen to agree with \(W\), since, by hypothesis, the numbers \(\lambda,\allowbreak \lambda',\allowbreak \lambda''\), etc., must be congruent to \(\lambda,\allowbreak \lambda h,\allowbreak \lambda hh\ldots\lambda h^{f-1}\) modulo \(n\) (the order is immaterial here), and therefore so must also \[\begin{gathered}
\lambda+\mu,\ \lambda'+\mu,\ \lambda''+\mu\quad\mathrm{etc.}\quad\text{to}\\
\lambda+\mu,\ \lambda h+\mu,\ \lambda hh+\mu\ldots\lambda h^{f-1}+\mu.\quad\textit{Q. E. D.}
\end{gathered}\]
\SourceLeaf{428}

We add the following corollaries to this theorem:

I. If \(k\) denotes any integer, the product of \((f,\allowbreak k\lambda)\) and \((f,\allowbreak k\mu)\) will be \[=(f,k(\lambda+\mu))+(f,k(\lambda^{\prime}+\mu))+(f,k(\lambda^{\prime\prime}+\mu))+\text{etc.}\]

II. Since each part of which \(W\) consists agrees either with the sum \((f,\allowbreak 0)\), which is \(=f\), or with one of \((f,\allowbreak 1),\allowbreak \ (f,\allowbreak g),\allowbreak \ (f,\allowbreak gg)\ldots(f,\allowbreak g^{e-1})\), \(W\) can be reduced to the following form \[\begin{aligned}
W=af+b(f,1)+b^{\prime}(f,g)&+b^{\prime\prime}(f,gg)\\
&+\ldots+b^\varepsilon(f,g^{e-1})
\end{aligned}\] where the coefficients \(a,\allowbreak b,\allowbreak b^{\prime}\), etc., will be positive integers (or some of them may also be \(=0\)). It is further clear that the product of \((f,\allowbreak k\lambda)\) and \((f,\allowbreak k\mu)\) then becomes \[=af+b(f,k)+b^{\prime}(f,kg)+\ldots+b^\varepsilon(f,kg^{e-1})\] Thus, \textit{for example}, for \(n=19\), the product of the sum \((6.1)\) with itself, or the square of this sum, becomes \(=(6,\allowbreak 2)+(6,\allowbreak 8)+(6,\allowbreak 9)+(6.12)+(6,\allowbreak 13)+(6,\allowbreak 19)=6+2(6,\allowbreak 1)+(6,\allowbreak 2)+2(6,\allowbreak 4).\)

III. Since the product of each part of \(W\) and a similar period \((f,\allowbreak \nu)\) can be reduced to an analogous form, it is clear that the product of three periods \((f,\allowbreak \lambda).(f,\allowbreak \mu).(f,\allowbreak \nu)\) can also be represented by \(cf+d(f,\allowbreak 1)\ldots+d^\varepsilon(f,\allowbreak g^{e-1})\), and that the coefficients \(c,\allowbreak d\), etc., will be integers and positive (or \(=0\)); moreover, for any integer value of \(k\), \[(f,k\lambda).(f,k\mu).(f,k\nu)=cf+d(f,k)+d^{\prime}(f,kg)+\text{etc.}\] In the same way this theorem extends to products of any number of similar periods, and it makes no difference whether these periods are all distinct, or partly or all identical.

IV. It follows that, if in any rational integral algebraic function \(F=\varphi(t,\allowbreak u,\allowbreak v\ldots)\) the indeterminates \(t,\allowbreak u,\allowbreak v\), etc., are replaced respectively by similar periods \((f,\allowbreak \lambda),\allowbreak (f,\allowbreak \mu),\allowbreak (f,\allowbreak \nu)\), etc., its value can be reduced to the form \[A+B(f,1)+B^{\prime}(f,g)+B^{\prime\prime}(f,gg)+\ldots+B^\varepsilon(f,g^{e-1})\] and the coefficients \(A,\allowbreak B,\allowbreak B^{\prime}\), etc., will all be integers if all the fixed coefficients in \(F\) are integers. If we then replace \(t,\allowbreak u,\allowbreak v\), etc., respectively
\SourceLeaf{429}
by \((f,\allowbreak k\lambda),\allowbreak (f,\allowbreak k\mu),\allowbreak (f,\allowbreak k\nu)\), etc., the value of \(F\) reduces to \(A+B(f,\allowbreak k)+B^{\prime}(f,\allowbreak kg)+\text{etc.}\)

\Needspace{6\baselineskip}\Article{346}
\textsc{Theorem.} \textit{Assuming that \(\lambda\) is a number not divisible by \(n\), and writing} \textit{\(p\) for \((f,\allowbreak \lambda)\) for brevity, any other similar period \((f,\allowbreak \mu)\), where \(\mu\) is also assumed not divisible by \(n\),} \textit{can be reduced to the following form} \[\alpha+\beta p+\gamma pp+\ldots+\theta p^{e-1}\] \textit{so that the coefficients \(\alpha,\allowbreak \beta\), etc., are fixed rational quantities.}

\textit{Proof.} For brevity, denote the periods \((f,\allowbreak \lambda g),\allowbreak (f,\allowbreak \lambda gg),\allowbreak (f,\allowbreak \lambda g^3)\), etc., up to \((f,\allowbreak \lambda g^{e-1})\), whose number is \(e-1\), and with one of which \((f,\allowbreak \mu)\) must agree, by \(p^{\prime},\allowbreak p^{\prime\prime},\allowbreak p^{\prime\prime\prime}\), etc. We then immediately have the equation \[0=1+p+p^{\prime}+p^{\prime\prime}+p^{\prime\prime\prime}+\text{etc.}\quad\mathinner{\ldotp\ \ldotp\ \ldotp\ \ldotp\ \ldotp}\quad(\mathrm I)\] Expanding the values of the powers of \(p\) up to the power with exponent \(e-1\), according to the rules of the preceding article, gives \(e-2\) further equations of the following kind. \[\begin{aligned}0&=pp+A+ap+a^{\prime} p^{\prime}+a^{\prime\prime} p^{\prime\prime}+a^{\prime\prime\prime} p^{\prime\prime\prime}+\text{etc.}\quad\mathinner{\ldotp\ \ldotp}\quad(\mathrm{II})\\0&=p^3+B+bp+b^{\prime} p^{\prime}+b^{\prime\prime} p^{\prime\prime}+b^{\prime\prime\prime} p^{\prime\prime\prime}+\text{etc.}\quad\mathinner{\ldotp\ \ldotp}\quad(\mathrm{III})\\0&=p^4+C+cp+c^{\prime} p^{\prime}+c^{\prime\prime} p^{\prime\prime}+c^{\prime\prime\prime} p^{\prime\prime\prime}+\text{etc.}\quad\mathinner{\ldotp\ \ldotp}\quad(\mathrm{IV})\quad\text{etc.}\end{aligned}\] Here all the coefficients \(A,\allowbreak a,\allowbreak a^{\prime}\), etc., \(B,\allowbreak b,\allowbreak b^{\prime}\), etc., etc., will be integers and, as should be carefully noted and follows directly from the preceding article, completely independent of \(\lambda\); \textit{that is}, the same equations will continue to hold whatever other value is assigned to \(\lambda\). This observation clearly extends to equation I as well, provided \(\lambda\) is taken not divisible by \(n\). --- Let us assume \((f,\allowbreak \mu)=p^{\prime}\); for it is readily seen that, if \((f,\allowbreak \mu)\) agrees with another of the periods \(p^{\prime\prime},\allowbreak p^{\prime\prime\prime}\), etc., arguments entirely analogous to those that follow can be used. Since the number of equations I, II, III, etc., is \(e-1\), the quantities \(p^{\prime\prime},\allowbreak p^{\prime\prime\prime}\), etc., whose number is \(=e-2\), can be eliminated from them by known methods, giving an equation of the following kind \((Z)\), free of those quantities, \[0=\mathfrak A+\mathfrak Bp+\mathfrak Cpp+\text{etc.}+\mathfrak Mp^{e-1}+\mathfrak Np^{\prime}\] This can be done so that all the coefficients \(\mathfrak A,\allowbreak \mathfrak B\ldots\mathfrak N\) are integers and certainly not all \(=0\). Now, if \(\mathfrak N=0\) does not hold, it is immediately clear that \(p^{\prime}\) is thereby determined in the way stated in the
\SourceLeaf{430}
theorem. It therefore remains to show that \(\mathfrak N=0\) cannot occur.

Assuming \(\mathfrak N=0\), equation \(Z\) becomes \(\mathfrak Mp^{e-1}+\text{etc.}+\mathfrak Bp+\mathfrak A=0\); since its degree certainly does not exceed \(e-1\), no more than \(e-1\) distinct values of \(p\) can satisfy it. But since the equations from which \(Z\) was derived are independent of \(\lambda\), it is clear that \(Z\) also does not depend on \(\lambda\), or holds whatever integer not divisible by \(n\) is taken for \(\lambda\). Thus equation \(Z\) will be satisfied whichever of the \(e\) sums \((f,\allowbreak 1),\allowbreak (f,\allowbreak g),\allowbreak (f,\allowbreak gg)\ldots(f,\allowbreak g^{e-1})\) is set equal to \(p\). It follows directly that these sums cannot all be unequal, but at least two must be equal to one another. Let one of two such equal sums contain the roots \([\zeta],\allowbreak [\zeta^{\prime}],\allowbreak [\zeta^{\prime\prime}]\), etc., and the other \([\eta],\allowbreak [\eta^{\prime}],\allowbreak [\eta^{\prime\prime}]\), etc.; assume (as we may) that all the numbers \(\zeta,\allowbreak \zeta^{\prime},\allowbreak \zeta^{\prime\prime}\), etc., \(\eta,\allowbreak \eta^{\prime},\allowbreak \eta^{\prime\prime}\), etc., are positive and \(<n\); clearly they will also all be distinct, and none \(=0\). Denote the function \[x^\zeta+x^{\zeta^{\prime}}+x^{\zeta^{\prime\prime}}+\text{etc.}-x^\eta-x^{\eta^{\prime}}-x^{\eta^{\prime\prime}}-\text{etc.}\] whose highest term does not exceed \(x^{n-1}\) by \(Y\). Clearly \(Y=0\) when \(x=[1]\) is set; hence \(Y\) contains the factor \(x-[1]\), which it will have in \textit{common} with the function denoted above by \(X\). But this is easily shown to be absurd. For if \(Y\) had any common factor with \(X\), the \textit{greatest} common divisor of the functions \(X,\allowbreak Y\) (which certainly cannot reach degree \(n-1\), as is already clear from the fact that \(Y\) is divisible by \(x\)) would have all its coefficients rational, as follows directly from the nature of the operations for finding the greatest common divisor of two such functions whose coefficients are all rational. But in Article 341 we showed that \(X\) cannot contain a factor of degree less than \(n-1\) whose coefficients are all rational. Therefore the assumption \(\mathfrak N=0\) cannot hold.

\textit{Example.} For \(n=19,\allowbreak \ f=6\), \(pp=6+2p+p^{\prime}+2p^{\prime\prime}\); from this and \(0=1+p+p^{\prime}+p^{\prime\prime}\) we deduce \(p^{\prime}=4-pp,\allowbreak \ p^{\prime\prime}=-5-p+pp\). Thus \[\begin{aligned}(6,2)&=4-(6,1)^2,\ &(6,4)&=-5-(6,1)+(6,1)^2\\(6,\allowbreak 4)&=4-(6,\allowbreak 2)^2,\allowbreak \ &(6.1)&=-5-(6,\allowbreak 2)+(6,\allowbreak 2)^2\\(6,\allowbreak 1)&=4-(6,\allowbreak 4)^2,\allowbreak \ &(6,\allowbreak 2)&=-5-(6,\allowbreak 4)+(6,\allowbreak 4)^2\end{aligned}\]
\SourceLeaf{431}


